\documentclass[12pt,a4paper]{article}
\usepackage[T1]{fontenc}
\usepackage{amsmath,amssymb,mathrsfs,tikz,times,pifont}
\usepackage{enumitem}
\usepackage{multicol}
\usepackage{lmodern}
\usetikzlibrary{trees}
\newcommand\circitem[1]{%
\tikz[baseline=(char.base)]{
\node[circle,draw=gray, fill=red!55,
minimum size=1.2em,inner sep=0] (char) {#1};}}
\newcommand\boxitem[1]{%
\tikz[baseline=(char.base)]{
\node[fill=cyan,
minimum size=1.2em,inner sep=0] (char) {#1};}}
\setlist[enumerate,1]{label=\protect\circitem{\arabic*}}
\setlist[enumerate,2]{label=\protect\boxitem{\alph*}}
%%%::::::by chnini ameur :::::::%%%
\everymath{\displaystyle}
\usepackage[left=1cm,right=1cm,top=1cm,bottom=1.7cm]{geometry}
\usepackage[colorlinks=true, linkcolor=blue, urlcolor=blue, citecolor=blue]{hyperref}
\usepackage{array,multirow}
\usepackage[most]{tcolorbox}
\usepackage{varwidth}
\usepackage{float} %pour utiliser l'option [H] qui force l'image à apparaître exactement à l'endroit où elle est placée dans le code.
\tcbuselibrary{skins,hooks}
\usetikzlibrary{patterns}
%%%::::::by chnini ameur :::::::%%%
\newtcolorbox{exa}[2][]{enhanced,breakable,before skip=2mm,after skip=5mm,
colback=yellow!20!white,colframe=black!20!blue,boxrule=0.5mm,
attach boxed title to top left ={xshift=0.6cm,yshift*=1mm-\tcboxedtitleheight},
fonttitle=\bfseries,
title={#2},#1,
% varwidth boxed title*=-3cm,
boxed title style={frame code={
\path[fill=tcbcolback!30!black]
([yshift=-1mm,xshift=-1mm]frame.north west)
arc[start angle=0,end angle=180,radius=1mm]
([yshift=-1mm,xshift=1mm]frame.north east)
arc[start angle=180,end angle=0,radius=1mm];
\path[left color=tcbcolback!60!black,right color = tcbcolback!60!black,
middle color = tcbcolback!80!black]
([xshift=-2mm]frame.north west) -- ([xshift=2mm]frame.north east)
[rounded corners=1mm]-- ([xshift=1mm,yshift=-1mm]frame.north east)
-- (frame.south east) -- (frame.south west)
-- ([xshift=-1mm,yshift=-1mm]frame.north west)
[sharp corners]-- cycle;
},interior engine=empty,
},interior style={top color=yellow!5}}
%%%%%%%%%%%%%%%%%%%%%%%
\usepackage{fancyhdr}
\usepackage{eso-pic}         % Pour ajouter des éléments en arrière-plan
% Commande pour ajouter du texte en arrière-plan
\usepackage{tkz-tab}
\AddToShipoutPicture{
    \AtTextCenter{%
        \makebox[0pt]{\rotatebox{80}{\textcolor[gray]{0.7}{\fontsize{5cm}{5cm}\selectfont PGB}}}
    }
}
\usepackage{lastpage}
\fancyhf{}
\pagestyle{fancy}
\renewcommand{\footrulewidth}{1pt}
\renewcommand{\headrulewidth}{0pt}
\renewcommand{\footruleskip}{10pt}
\fancyfoot[R]{
\color{blue}\ding{45}\ \textbf{2026}
}
\fancyfoot[L]{
\color{blue}\ding{45}\ \textbf{Prof:M. BA}
}
\cfoot{\bf
\thepage /
\pageref{LastPage}}
% Création du compteur pour les exercices
\newcounter{exercice}
\renewcommand{\theexercice}{\arabic{exercice}}  % Définit l'affichage du compteur en chiffres arabes

% Définir la commande \exo
\newcommand{\exo}{\refstepcounter{exercice}\textbf{Exercice \theexercice} }
\begin{document}
\renewcommand{\arraystretch}{1.5}
\renewcommand{\arrayrulewidth}{1.2pt}
\begin{tikzpicture}[overlay,remember picture]
    \node[draw=blue,line width=1.2pt,fill=purple,text=blue,inner sep=3mm,rounded corners,pattern=dots]at ([yshift=-2.5cm]current page.north) {\begingroup\setlength{\fboxsep}{0pt}\colorbox{white}{\begin{tabular}{|*1{>{\centering \arraybackslash}p{0.28\textwidth}} |*2{>{\centering \arraybackslash}p{0.2\textwidth}|} *1{>{\centering \arraybackslash}p{0.19\textwidth}|} }
                \hline
                \multicolumn{3}{|c|}{$\diamond$$\diamond$$\diamond$\ \textbf{Lycée de Dindéfélo}\ $\diamond$$\diamond$$\diamond$ } & \textbf{A.S. : 2026/2027}                                              \\ \hline
                \textbf{Matière: Mathématiques}                                                                                    & \textbf{Niveau : 1}\textbf{L} & \textbf{Date: 16/11/2026} & \textbf{} \\ \hline
                \multicolumn{4}{|c|}{\parbox[c]{10cm}{\begin{center}
                                                                  \textbf{{\Large\sffamily Td Polynomes Correction}}
                                                              \end{center}}}                                                                                                        \\ \hline
            \end{tabular}}\endgroup};
\end{tikzpicture}

\vspace{3cm}

\fbox{\textbf{Correction 1}}

% --- IMAGE 6 : POLYNÔMES ---
\textbf{A) Compléter les phrases suivantes}
\begin{enumerate}
    \item Si $\alpha$ est racine d'un polynôme $P(x)$ alors $P(x)$ est factorisable par $\mathbf{(x - \alpha)}$.
    \item La forme factorisée d'un trinôme $ax^2 + bx + c$ dont les racines sont $x_1$ et $x_2$ est $\mathbf{a(x - x_1)(x - x_2)}$.
    \item Si $\alpha$ est une racine d'un polynôme de degré 4. Il existe $Q(x)$ de degré $\mathbf{3}$ tel que $P(x) = (x - \alpha)Q(x)$. \\
    \textit{Remarque :} Si le coefficient du plus haut degré de $P(x)$ est $a$, alors le coefficient du plus haut degré de $Q(x)$ est également $\mathbf{a}$ (on peut aussi écrire $P(x) = a(x - \alpha)R(x)$ où $R(x)$ est un polynôme unitaire de degré $3$).
\end{enumerate}

\textbf{B) Répondre par vrai ou faux}
\begin{enumerate}
    \item $x^3 - y^3 = (x + y)(x^2 - xy + y^2)$ : \textbf{Faux.} \\
    \textit{Justification :} La formule correcte est $x^3 - y^3 = (x - y)(x^2 + xy + y^2)$.
    
    \item Un polynôme de degré 3 admet 3 racines : \textbf{Faux.} \\
    \textit{Justification :} Dans $\mathbb{R}$, un polynôme de degré 3 admet \textbf{au plus} 3 racines (par exemple, $P(x) = x^3 + x$ n'admet que $0$ comme racine réelle). \\
    \textit{Cas général (si $P(x)$ admet 3 racines distinctes $x_1, x_2, x_3$) :} Si son coefficient dominant est $a$, sa forme factorisée est $P(x) = a(x - x_1)(x - x_2)(x - x_3)$.
    
    \item Un polynôme qui a pour racine $\alpha$ et $\beta$ est factorisable par $(x - \alpha)(x - \beta)$ : \textbf{Faux} (en général). \\
    \textit{Justification :} Cela n'est vrai que si $\alpha$ et $\beta$ sont \textbf{distincts} ($\alpha \neq \beta$). Dans ce cas, il existe un polynôme $Q(x)$ tel que $P(x) = (x - \alpha)(x - \beta)Q(x)$ (où $Q(x)$ a le même coefficient dominant $a$ que $P(x)$ ; en particulier, si $P$ est de degré $2$, alors $Q(x) = a$ et $P(x) = a(x - \alpha)(x - \beta)$).
\end{enumerate}

\fbox{\textbf{Correction 2}}

Complète les pointillés

\begin{enumerate}
    \item On appelle racine réelle d'un polynôme $P$ tout nombre réel $\alpha$ tel que $\mathbf{P(\alpha) = 0}$.
    \item Si un polynôme $P$ a une racine réelle $\alpha$, alors il existe un polynôme $Q$ tel que $P(x) = \mathbf{(x - \alpha)Q(x)}$.
    \item Soit $P(x) = ax^2 + bx + c$ un trinôme du second degré.
    \begin{enumerate}
        \item Si $P(x)$ admet deux racines $x_1$ et $x_2$ distinctes, alors sa forme factorisée est $P(x) = \mathbf{a(x - x_1)(x - x_2)}$.
        \item Si $P(x)$ admet une racine double $x_0$, alors sa forme factorisée est $P(x) = \mathbf{a(x - x_0)^2}$.
    \end{enumerate}
\end{enumerate}

\fbox{\textbf{Correction 3}}

Soit le polynôme $P$ défini par : $P(x) = 2x^3 - 9x^2 - 8x + 15$.

\begin{enumerate}
    \item \textbf{Calcul de $P(5)$ :}
    \begin{align*}
        P(5) &= 2(5)^3 - 9(5)^2 - 8(5) + 15 \\
             &= 2 \times 125 - 9 \times 25 - 40 + 15 \\
             &= 250 - 225 - 40 + 15 \\
             &= 25 - 40 + 15 \\
             &= 0
    \end{align*}
    Donc $5$ est une racine de $P$.

    \item \textbf{Factorisation de $P(x)$ :} \\
    Puisque $P(5) = 0$ et que $\deg(P) = 3$, il existe un polynôme $Q(x)$ de degré $2$ tel que $P(x) = (x - 5)Q(x)$.

    \begin{enumerate}
        \item \textbf{Méthode d'identification des coefficients :} \\
        Posons $Q(x) = ax^2 + bx + c$ (avec $a \neq 0$). Développons l'expression $(x - 5)(ax^2 + bx + c)$ :
        \begin{align*}
            (x - 5)(ax^2 + bx + c) &= ax^3 + bx^2 + cx - 5ax^2 - 5bx - 5c \\
                                   &= ax^3 + (b - 5a)x^2 + (c - 5b)x - 5c
        \end{align*}
        Par identification avec $P(x) = 2x^3 - 9x^2 - 8x + 15$, on obtient le système :
        \[
        \begin{cases}
            a = 2 \\
            b - 5a = -9 \\
            c - 5b = -8 \\
            -5c = 15
        \end{cases}
        \iff
        \begin{cases}
            a = 2 \\
            b = -9 + 5(2) = 1 \\
            c = \dfrac{15}{-5} = -3 \\
            c - 5b = -3 - 5(1) = -8 \quad \text{(vérifié)}
        \end{cases}
        \]
        Ainsi, $Q(x) = 2x^2 + x - 3$ et $P(x) = (x - 5)(2x^2 + x - 3)$.

        \item \textbf{Méthode de la division euclidienne :} \\
        En effectuant la division euclidienne de $2x^3 - 9x^2 - 8x + 15$ par $(x - 5)$ :
        \begin{itemize}
            \item $2x^3 \div x = 2x^2$ ; on soustrait $2x^2(x - 5) = 2x^3 - 10x^2$, il reste $x^2 - 8x + 15$.
            \item $x^2 \div x = x$ ; on soustrait $x(x - 5) = x^2 - 5x$, il reste $-3x + 15$.
            \item $-3x \div x = -3$ ; on soustrait $-3(x - 5) = -3x + 15$, il reste $0$.
        \end{itemize}
        On retrouve le quotient $Q(x) = 2x^2 + x - 3$, d'où :
        \[
            P(x) = (x - 5)(2x^2 + x - 3)
        \]

        \textbf{Factorisation complète de $P(x)$ :} \\
        Factorisons le trinôme $Q(x) = 2x^2 + x - 3$. On remarque que $1$ est une racine évidente ($2(1)^2 + 1 - 3 = 0$), ou par le discriminant $\Delta = 1^2 - 4(2)(-3) = 25 > 0$ :
        \[
            x_1 = \frac{-1 - 5}{4} = -\frac{3}{2} \quad \text{et} \quad x_2 = \frac{-1 + 5}{4} = 1
        \]
        Donc $2x^2 + x - 3 = 2(x - 1)\left(x + \frac{3}{2}\right) = (x - 1)(2x + 3)$. \\
        Finalement, la forme complètement factorisée de $P(x)$ est :
        \[
            \mathbf{P(x) = (x - 5)(x - 1)(2x + 3)}
        \]
    \end{enumerate}
\end{enumerate}
\fbox{\textbf{\exo}}
Soit $P$ le polynôme défini par $P(x) = 3x^3 - x^2 - 8x - 4$.

\begin{enumerate}
    \item Vérifier que $2$ est une racine de $P$.
    \item Déterminer le polynôme $Q$ tel que $P(x) = (x - 2)Q(x)$.
    \item Factoriser $Q(x)$ puis en déduire une factorisation complète de $P(x)$.
    \item Résoudre dans $\mathbb{R}$ :
        \begin{enumerate}
        \item $P(x) = 0$ 
        \item $P(x + 2) = 0$ 
        \item $P(x) < 0$
        \end{enumerate}
\end{enumerate}

\fbox{\textbf{\exo}}
Soit les polynômes $P(x) = -3x^3 - 2x^2 + 19x - 6$ et $T(x) = (1 - 3x)(2 - x)(-x - 3)$.

\begin{enumerate}
    \item 
    \begin{enumerate} 
    \item Calculer $P(1)$, $P(-3)$ et $P(0)$.
    
    \item Lequel des réels $1$, $-3$ et $0$ est une racine de $P$ ? Justifie
    \end{enumerate}
    \item Écrire $P(x)$ sous forme de produit de facteurs du premier degré. 
    
    \item 
    \begin{enumerate} 
    \item Donne en justifiant le degré du polynôme $T$. 
    \item Donner le monôme de plus haut degré de $T$. 
    
    \item Résoudre dans $\mathbb{R}$ l'équation $P(x) = 0$.
    
    \item Résoudre dans $\mathbb{R}$ l'inéquation $P(x) \leq 0$.
    \end{enumerate}
    \item Montrer que les polynômes $P$ et $T$ sont égaux.
\end{enumerate}

\fbox{\textbf{\exo}}

Soit le polynôme $f(x) = x^3 + ax^2 + bx + 6$ avec $a$ et $b$ sont deux nombres réels.

\begin{enumerate}
    \item Déterminer $a$ et $b$ sachant que $f(-2) = 0$ et $f(-1) = 8$.
    \item On pose $f(x) = x^3 - 2x^2 - 5x + 6$.
    \begin{enumerate}
        \item Calculer $f(1)$. Que peut-on en déduire ?
        \item Déterminer un polynôme $Q(x)$ tel que $f(x) = (x - 1)Q(x)$.
        \item Factoriser $Q(x)$ et en déduire une factorisation (complètement) de $f(x)$.
        \item Résoudre dans $\mathbb{R}$ :
        \begin{enumerate}
            \item $f(x) = 0$
            \item $f(x) \geq 0$
        \end{enumerate}
    \end{enumerate}
\end{enumerate}

\fbox{\textbf{\exo}}

\begin{enumerate}
    \item[I.] Définir les termes suivants : \\
    \textbf{Trinôme du second - Monôme de la variable $x$ - Polynôme}
    
    \item[II.] Résoudre dans $\mathbb{R}$ les équations et inéquations suivantes :
\begin{multicols}{2}
\setlength{\columnseprule}{0.1mm} % La largeur de la ligne verticale entre les colonnes
    \begin{enumerate}
        \item $x^2 - 6x = 0$
        \item $-2x^2 + 3x + 2 = 0$
        \item $x^2 - 5x + 4 \leq 0$
        \item $-x^2 - x - 6 > 0$ 
    \end{enumerate}
\end{multicols}
\end{enumerate}

\fbox{\textbf{\exo}}

 On considère le polynôme $P$ défini par : $P(x) = x^3 - x^2 - 14x + 24$

\begin{enumerate}
    \item Calculer $P(1)$ ; $P(2)$ et $P(-3)$. 
    \item En déduire une racine de $P$. 
    \item Montrer qu'il existe un polynôme $Q$ tel que : $P(x) = (x - 2)Q(x)$ 
    \item Déterminer $Q(x)$. En déduire une factorisation complète de $P(x)$ 
    \item Résoudre dans $\mathbb{R}$ l'équation : $P(x) = 0$. En déduire les solutions dans $\mathbb{R}$ : \\
    $P(|x|) = 0$ 
    \vspace{0.3cm}
    \item Résoudre dans $\mathbb{R}$ l'équation : $P(x) \leq 0$ 
\end{enumerate}

\fbox{\textbf{\exo}}

\begin{enumerate}
    \item Trouver a et b tels que $2$ et $-1$ sont des racines de $P(x) = ax^3 + bx^2 - 5x - a$.
    
    \item En prenant les valeurs de a et b trouver dans 1), factoriser complètement $P(x)$.
    
    \item Résoudre dans $\mathbb{R}$ l'équation $P(x) = 0$. 
    
    \item En déduire les solutions des équations :
    \begin{enumerate}
        \item $P(x - 1) = 0$.
        \item $P(-|x|) = 0$. 
    \end{enumerate}
    
    \item Résoudre dans $\mathbb{R}$ l'inéquation $P(x) < 0$.
\end{enumerate}

\fbox{\textbf{\exo}} On considère le polynôme $P(x) = x^3 - 2x^2 - x + 2$
 \begin{multicols}{2}
\setlength{\columnseprule}{0.1mm} % La largeur de la ligne verticale entre les colonnes  
\begin{enumerate}
    \item Montrer que 1 est racine de $P(x)$. 
    
    \item En déduire une factorisation de $P(x)$.
    
    \item Résoudre l'équation $P(x) = 0$.
    
    \item Déduire les solutions des équations :
    \begin{enumerate}
        \item $(-x + 1)^3 - 2(-x + 1)^2 - (-x + 1) + 2 = 0$.
        \item $\dfrac{1}{(x+3)^3} - \dfrac{2}{(x+3)^2} - \dfrac{1}{x+3} + 2 = 0$.
    \end{enumerate}
    \item Résoudre l'inéquation $\dfrac{x^3 - 2x^2 - x + 2}{x^2 - x - 6} > 0$.
\end{enumerate}
 \end{multicols}

  \begin{multicols}{2}
\setlength{\columnseprule}{0.1mm} % La largeur de la ligne verticale entre les colonnes 
\fbox{\textbf{\exo}}
On donne $P(x) = 2x^4 - 7x^3 - 2x^2 + 13x + 6$

\begin{enumerate}
    \item Calculer $P(0)$ et $P(-1)$ puis conclure. 
    \item Déterminer le polynôme $Q(x)$ tel que $P(x) = (x + 1)Q(x)$ 
    \item Dans la suite de l'exercice, on suppose que $Q(x) = 2x^3 - 9x^2 + 7x + 6$    
    \begin{enumerate}
        \item Sachant que 2 est une racine de $Q(x)$, factoriser complètement $Q(x)$. 
        \item En déduire une factorisation de $P(x)$.
        \item Résoudre dans $\mathbb{R}$ l'équation $P(x) = 0$.
        \item Déduire de la question 3c) la résolution dans $\mathbb{R}$ de l'équation $P(x + 1) = 0$.
        \item Résoudre dans $\mathbb{R}$ l'inéquation $P(x) \leq 0$.
    \end{enumerate}
\item On pose : $f(x) = \dfrac{Q(x)}{-x^2 + 7x - 10}$
    \begin{enumerate}
        \item Déterminer le domaine de définition de $f$.
        \item Simplifier $f(x)$ puis résoudre dans $\mathbb{R}$ l'inéquation $f(x) > 0$.
    \end{enumerate}
\end{enumerate}
\end{multicols}
\fbox{\textbf{\exo}}

On considère le polynôme $P$ défini par : $P(x) = x^3 + 3x^2 - 4$.
 \begin{multicols}{2}
\setlength{\columnseprule}{0.1mm} % La largeur de la ligne verticale entre les colonnes  
\begin{enumerate}
    \item Montrer que 1 est une racine de $P$.
    \item Donner une factorisation de $P(x)$.
    \item Résoudre dans $\mathbb{R}$ : 
    \begin{enumerate}
        \item l'équation $P(x) = 0$.
        \item l'inéquation $P(x) \geq 0$.
    \end{enumerate}
\end{enumerate}
\end{multicols}

\fbox{\textbf{\exo}}

On considère le polynôme $P$ défini par : $P(x) = -x^3 + 6x^2 - 9x + k$ où $k$ est un nombre réel.

\begin{enumerate}
    \item Déterminer la valeur du réel $k$ pour que 4 soit une racine de $P$.
    \item Pour la valeur de $k$ trouvée en 1) :
\begin{multicols}{2}
\setlength{\columnseprule}{0.1mm} % La largeur de la ligne verticale entre les colonnes
    \begin{enumerate}
        \item Factoriser $P$.
        \item Résoudre dans $\mathbb{R}$ l'équation $P(x) = 0$.
        \item Résoudre dans $\mathbb{R}$ l'inéquation $P(x) < 0$.
    \end{enumerate}
\end{multicols}
\end{enumerate}

\end{document}